%\documentclass[draft, 11pt, a4paper]{article}  %compiles faster for editing
\documentclass[11pt, a4paper]{article}
%\usepackage{geometry}
\usepackage[inner=1.5cm,outer=1.5cm,top=2.5cm,bottom=2.5cm]{geometry}
\pagestyle{empty}
\usepackage{graphicx}
\usepackage{fancyhdr, lastpage, bbding, pmboxdraw}
\usepackage[usenames,dvipsnames]{color}
\usepackage{epigraph}

\newcommand{\hilight}[1]{\colorbox{BurntOrange}{#1}}  % changes will be highlighted in BurntOrange

% Set links, etc. to different colors
\definecolor{OliveGreen}{cmyk}{0.64,0,0.95,0.40}
\definecolor{CadetBlue}{cmyk}{0.62,0.57,0.23,0}
\definecolor{lightlightgray}{gray}{0.93}
\definecolor{darkblue}{rgb}{0,0,.6}
\definecolor{darkred}{rgb}{.7,0,0}
\definecolor{darkgreen}{rgb}{0,.6,0}
\definecolor{red}{rgb}{.98,0,0}
\usepackage[bookmarksopen, backref,colorlinks,pagebackref,pdfusetitle,urlcolor=darkblue,citecolor=darkblue,linkcolor=darkred,bookmarksnumbered,plainpages=false]{hyperref}


\renewcommand{\thefootnote}{\fnsymbol{footnote}}

\pagestyle{fancyplain}
\fancyhf{}
\lhead{ \fancyplain{}{Product Design Specifications} }
%\chead{ \fancyplain{}{} }
%\rhead{ \fancyplain{}{} }
%\rfoot{\fancyplain{}{page \thepage\ of \pageref{LastPage}}}
\fancyfoot[R] {\thepage}     %use this instead
\thispagestyle{plain}

%%%%%%%%%%%% LISTING %%%
\usepackage{listings}
\usepackage{caption}
\DeclareCaptionFont{white}{\color{white}}
\DeclareCaptionFormat{listing}{\colorbox{gray}{\parbox{\textwidth}{#1#2#3}}}
\captionsetup[lstlisting]{format=listing,labelfont=white,textfont=white}
\usepackage{fancyvrb}
\usepackage{acronym}
\usepackage{amsthm}

%%% My packages:
\usepackage{graphicx} %package to manage images
\graphicspath{ {images/} }  %set default image path
\usepackage{chngcntr}
\usepackage{float}       % useful for positioning figures / etc.
\counterwithin{table}{section}      % counts fig 1.1... 1.2... etc
\counterwithin{figure}{section}     % same for figures.. for List of Figures page

% for highlighting text
\usepackage{soul}
\newcommand{\hlc}[2][yellow]{ {\sethlcolor{#1} \hl{#2}} }  %Not sure why I set it to yellow.. maybe used for something else besides changes
\usepackage[export]{adjustbox} %for left/right aligned stuff
\usepackage{array} % for centering things in tables
\newcolumntype{M}[1]{>{\centering\arraybackslash}m{#1}} %SAME AS ABOVE

\lstset{
%language=bash,                          % Code language
basicstyle=\ttfamily,                   % Code font, Examples: \footnotesize, \ttfamily
keywordstyle=\color{OliveGreen},        % Keywords font ('*' = uppercase)
commentstyle=\color{gray},              % Comments font
numbers=left,                           % Line nums position
numberstyle=\tiny,                      % Line-numbers fonts
stepnumber=1,                           % Step between two line-numbers
numbersep=5pt,                          % How far are line-numbers from code
backgroundcolor=\color{lightlightgray}, % Choose background color
frame=none,                             % A frame around the code
tabsize=2,                              % Default tab size
captionpos=t,                           % Caption-position = bottom
breaklines=true,                        % Automatic line breaking?
breakatwhitespace=false,                % Automatic breaks only at whitespace?
showspaces=false,                       % Dont make spaces visible
showtabs=false,                         % Dont make tabls visible
columns=flexible,                       % Column format
morekeywords={__global__, __device__},  % CUDA specific keywords
}

%%%%%%%%%%%%%%%%%%%%%%%%%%%%%%%%%%%%

\makeatletter                                   %fix for figure names and numbers overlapping
\def\l@figure{\@dottedtocline{1}{1.5em}{4.5em}}   % in the List Of Figures
\makeatother



\fancyfoot[L]{\small{\textit{Revision 1.1}}} %revision number here starting on TOC

\begin{document}

\pagenumbering{Roman} %roman numerals for title.. TOC.. list of figures/tables..
\thispagestyle{empty}   % removes page numbering for title page

\begin{figure}[t]
    \begin{center}
    {\LARGE \textsc{ Product Design Specifications}}
    \end{center}

    \begin{center}
    \Large Inductive Charging
    \end{center}

    \begin{center}
    \normalsize \textit{Practicum Team 10}
    \end{center}
\end{figure}

\vspace*{\fill}

\begingroup
    \begin{center}
    \includegraphics[scale=3.5]{MCECS.jpg}
    \end{center}
\endgroup
    
\vspace*{\fill}
    
\begin{figure}[b]
    \includegraphics[width=.3\textwidth,center]{pdx.png}
\end{figure}


% Contact info footer
    \begin{figure}[!b]
    \begin{center}
    \rule{6in}{0.4pt}
    \begin{minipage}[t]{.75\textwidth}
    \begin{tabular}{llcccll}
    \textbf{Team Members:}  Tyler Gilbert, Kati Dahn, Jann Messer, Jeff Brown & & & &\\
    \textbf{Project repository:} \href{https://github.com/tjgilbert/ece411}{https://github.com/tjgilbert/ece411} \\
    \end{tabular}
    \end{minipage}
    \rule{6in}{0.4pt}
    \end{center}
    \end{figure}
    
\cleardoublepage
\phantomsection

\hypertarget{MyToc}{} %For linking back to the TOC
\rhead{ \fancyplain{}{Table of Contents} }
\begin{center}
{\Large \textsc{ Table of Contents}}
\end{center}
\renewcommand{\contentsname}{}
\tableofcontents

\cleardoublepage
\pagenumbering{arabic}   %  Then switch to arabic page numbering at 1. 
\fancyfoot[R] {page \thepage\ of \pageref*{LastPage} }  



\rhead{ \fancyplain{}{Introduction} }
\section{Introduction}

\section{Project Overview}

\section{Requirements}

    \subsection{Functional Requirements}
        \begin{centering}
        \begin{tabular}{|l|m{7cm}|l|m{7cm}|} \hline
        \textit{\textbf{Req\#}}	& \textit{\textbf{Requirement}} &\textit{\textbf{Weight}}&\textit{\textbf{Justification}} \\ \hline
        1.1	& Have one or more sensors and actuators & Must	& Course requirement \\ \hline
        1.2	& Have a digital or analog processor & Must	& Course requirement \\ \hline
        1.3	& Wirelessly charge a USB device & Must	& This is the premise of the project, this is it's sole purpose \\ \hline
        1.4	& Turn on transmitter when device to be charged is placed in range	& Should & This will save energy because the product will not transmit energy if no device is present \\ \hline
        1.5 & Turn off transmitter when device is removed	& Should & Same as 1.4 \\ \hline
        \end{tabular}
        \end{centering}
    
    \subsection{Performance Requirements}
        \begin{centering}
        \begin{tabular}{|l|m{7cm}|l|m{7cm}|} \hline
        \textit{\textbf{Req\#}}	& \textit{\textbf{Requirement}} &\textit{\textbf{Weight}}&\textit{\textbf{Justification}} \\ \hline
        2.1	& Provide USB spec power (4.40-5.25V, 500mA)	& Should & Meeting this spec would guarantee it would work on most mobile devices \\ \hline
        2.2	& Can fit charging unit into a backpack (300mm x 300mm x 150mm)	& Should & Meeting these dimensions would qualify the charging unit as portable \\ \hline
        2.3	& Receiver is smaller or equal to 138mm x 67mm	& Should & This is the dimension of the iPhone 7, meeting this footprint would make it compatible with more devices \\ \hline
        \end{tabular}
        \end{centering}
    
    \subsection{Economic Requirements}
        \begin{centering}
        \begin{tabular}{|l|m{7cm}|l|m{7cm}|} \hline
        \textit{\textbf{Req\#}}	& \textit{\textbf{Requirement}} &\textit{\textbf{Weight}}&\textit{\textbf{Justification}} \\ \hline
        3.1 & Prototypes not to exceed \$75 each & Should & We are not made of money  \\ \hline
        3.2 & Product components not to exceed \$20 ea for quantity 1000 & Should & We are not paying our engineers, or technicians, If we sell the product for \$20+S\&H we \textit{might} turn a profit \\ \hline
        \end{tabular}
        \end{centering}    

    \subsection{Health and Safety Requirements}
        \begin{centering}
        \begin{tabular}{|l|m{7cm}|l|m{7cm}|} \hline
        \textit{\textbf{Req\#}}	& \textit{\textbf{Requirement}} &\textit{\textbf{Weight}}&\textit{\textbf{Justification}} \\ \hline
        4.1	& Shut off if base draws excessive current (current rating of Tx coil wire)	& Must &This will serve as short circuit protection, which could lead to damage of product, and/or fire \\ \hline
        4.2	& Shut off if the base unit overheats \mbox{(exceeds $35^{\circ}C$)}& Should & $35^{\circ}C$ is the maximum operating temperature of the iPhone \\ \hline
        \end{tabular}
        \end{centering}
    
    \subsection{Legal Requirements}
        \begin{centering}
        \begin{tabular}{|l|m{7cm}|l|m{7cm}|} \hline
        \textit{\textbf{Req\#}}	& \textit{\textbf{Requirement}} &\textit{\textbf{Weight}}&\textit{\textbf{Justification}} \\ \hline
        6.1	& Code and hardware design will be open source or generated by us& Must	&To preserve our academic honesty and integrity \\ \hline
        6.2	& Any borrowed or modified code/designs will be credited according to the applicable license & Must	&It's the right thing to do\\ \hline
        \end{tabular}
        \end{centering}
    
    \subsection{Environmental Requirements}
        \begin{centering}
        \begin{tabular}{|l|m{7cm}|l|m{7cm}|} \hline
        \textit{\textbf{Req\#}}	& \textit{\textbf{Requirement}} &\textit{\textbf{Weight}}&\textit{\textbf{Justification}} \\ \hline
        7.1	& Charging unit compliant with IP22	& Should & Users may spill liquids on the base unit, but as most mobile devices are not water resistant the receive units will not carry this rating \\ \hline
        \end{tabular}
        \end{centering}
    
    
    \subsection{Usability Requirements}
        \begin{centering}
        \begin{tabular}{|l|m{7cm}|l|m{7cm}|}        \hline
        \textit{\textbf{Req\#}}	& \textit{\textbf{Requirement}} &\textit{\textbf{Weight}}&\textit{\textbf{Justification}} \\ \hline
        8.1 & Orientation/location for device placement on charging unit is marked & Should & Marking the location for optimal charging will reduce the time to full charge\\ \hline
        8.2 & User able set up within 5 minutes without user manual by a college student & Should & This is a metric to decide if the device is user friendly \\ \hline
        8.3 & Base station will indicate when device is coupled/charging & Should & User feedback will increase confidence that the product is operating correctly                                \\ \hline
        8.4 & Base station will indicate when device is done charging & May& Further user feedback on charge state would would be a more obvious indication than having to check the phone \\ \hline
        \end{tabular}
        \end{centering}

\newpage   




\begin{thebibliography}{}

\bibitem{USB20}
"Universal Serial Bus Specification Revision 2.0"
\emph{Compaq Computer Corporation,
Hewlett-Packard Company, Intel Corporation, Lucent Technologies Inc,
Microsoft Corporation, NEC Corporation, Koninklijke Philips Electronics N.V.} \small{[\href{http://www.pjrc.com/teensy/beta/usb20.pdf\#page=206}{PDF DOCUMENT}]}, \normalsize 2000
 
\bibitem{appletemp}
"Keeping IPhone, IPad, and IPod Touch within Acceptable Operating Temperatures."
\emph{Apple Support.} N.p., 28 Sept. 2016. Web. 24 Oct. 2016. \url{https://support.apple.com/en-us/HT201678}

\bibitem{iprating}
"IP Rated Enclosures Explained"
\emph{The Enclosure Company (International) Ltd} N.p., n.d. Web. 23 Oct. 2016. \url{http://www.enclosurecompany.com/ip-ratings-explained.php}


\rhead{ \fancyplain{}{References} }

\end{thebibliography}
% that's all folks







    

    


    

    
    

    
    



    

    

    

    
    

\end{document}
	
	